\documentclass[a4paper,10pt]{article}
\usepackage[margin=18mm,footskip=9mm]{geometry}
\usepackage{iftex}
\usepackage{fontspec}
\setmainfont{texgyreheros}[Extension=.otf,UprightFont=*-regular,BoldFont=*-bold,ItalicFont=*-italic,BoldItalicFont=*-bolditalic]
\newif\ifChineseMemo
\ChineseMemotrue
\usepackage[UTF8,fontset=fandol]{ctex}
\renewcommand{\familydefault}{\sfdefault}
\usepackage{booktabs,graphicx,xcolor,caption,xurl,fancyhdr}
\usepackage{pgfplots}
\pgfplotsset{compat=1.18}
\usepgfplotslibrary{dateplot}
\definecolor{MemoBlue}{HTML}{203E70}
\definecolor{MemoRule}{HTML}{B9C8DF}
\usepackage[colorlinks=true,linkcolor=MemoBlue,urlcolor=MemoBlue]{hyperref}
\captionsetup{font=small,labelfont=bf,skip=4pt}
\setlength{\parindent}{0pt}
\setlength{\parskip}{5pt}
\pagestyle{fancy}
\fancyhf{}
\renewcommand{\headrulewidth}{0pt}
\renewcommand{\footrulewidth}{0.3pt}
\newcommand{\ReportTitle}{20日反转因子：样本外检验未通过}
\newcommand{\ReportType}{因子研究报告}
\newcommand{\ReportDate}{2026年11月15日}
\newcommand{\DataDate}{2026-11-13}
\newcommand{\ReportScope}{2023-11-16—2026-11-13；月度调仓；人民币计价}
\newcommand{\TeamName}{量化研究团队}
\newcommand{\DeskName}{期末项目 · 横截面因子}
\newcommand{\ReportAuthor}{生成组、检查组、分配组}
\newcommand{\Coverage}{34只创业板股票；35个可检验月份}
\newcommand{\KeyFinding}{样本内选出的20日反转因子在最后12个月失去预测能力，多空策略年化收益为-29.20\%，不应据此进入实盘。}
\fancyfoot[L]{\scriptsize\TeamName}
\fancyfoot[R]{\scriptsize\ReportType\quad\thepage}

\begin{document}
\begin{minipage}[t]{0.70\textwidth}
\vspace{0pt}
{\small\ReportDate}\par\medskip
{\color{MemoBlue}\hrule height 0.5pt}\smallskip
{\large\ReportType}\par
{\LARGE\bfseries\ReportTitle}\par\medskip
\ReportScope
\end{minipage}\hfill
\begin{minipage}[t]{0.23\textwidth}
\vspace{0pt}\raggedleft
\colorbox{MemoBlue}{\parbox[c][19mm][c]{0.86\linewidth}{\centering\color{white}\bfseries\TeamName}}\par\medskip
\DeskName
\end{minipage}
\par\bigskip
{\color{MemoBlue}\hrule height 0.5pt}\medskip
{\large\KeyFinding}\par\bigskip

\begin{minipage}[t]{0.61\textwidth}
\vspace{0pt}\raggedright
\setlength{\parskip}{4pt}
\textbf{研究问题}\par
在现有34只股票中，月末计算的因子能否预测下个月横截面收益？生成组仅比较动量、反转和低波动三个候选，分配组按前23个月的月均Rank IC选择因子，检查组再核验最后12个月。

\textbf{主要发现}\par
20日反转以样本内月均Rank IC 0.0149入选，但ICIR仅0.0640，信号很弱。样本外Rank IC降至-0.0048，方向反转；多空组合12个月累计及年化收益均为-29.20\%，最大回撤-49.77\%。同期等权基准收益为-1.80\%。

\textbf{解释}\par
反转因子偏好过去20日跌幅较大的股票。样本外Q5年化收益5.97\%，低于Q1的16.35\%，未形成预期的高因子值占优。样本内正Rank IC没有转化为可接受的收益或样本外稳定性，因此结论是“研究完成但未通过”，而不是“发现可交易alpha”。

\textbf{风险与限制}\par
数据仅覆盖34只创业板股票，并非完整沪深300；使用当前成分名单，存在幸存者偏差。收盘价复权口径未确认，未计分红、手续费、滑点、涨跌停和做空约束。12个月样本外期很短，结果不构成投资建议。
\end{minipage}\hfill
\begin{minipage}[t]{0.36\textwidth}
\vspace{0pt}\small
\captionof{table}{20日反转因子的核心检验结果}\label{tab:core-metrics}
{\footnotesize 月末信号、次月收益；样本外为最后12个月。\par}
{\scriptsize\begin{tabular}{@{}lrr@{}}
\toprule
指标 & 样本内 & 样本外 \\
\midrule
月数 & 23 & 12 \\
月均Rank IC & 0.0149 & -0.0048 \\
ICIR & 0.0640 & -0.0134 \\
多空年化收益 & -46.36\% & -29.20\% \\
多空年化波动 & 42.43\% & 62.57\% \\
Sharpe（零利率） & -1.20 & -0.28 \\
多空最大回撤 & -70.80\% & -49.77\% \\
\bottomrule
\end{tabular}
}

\medskip
\captionof{figure}{样本外累计净值}\label{fig:oos-wealth}
{\footnotesize 多空组合为Q5减Q1，逐月复利；基准为34只股票等权。\par}
\begin{tikzpicture}
\begin{axis}[width=\linewidth,height=55mm,date coordinates in=x,xticklabel={\year-\month},xticklabel style={rotate=35,anchor=east,font=\scriptsize},ylabel={净值（起点=100）},grid=major,grid style={gray!20},legend style={at={(0.02,0.98)},anchor=north west,font=\scriptsize},every axis label/.style={font=\small}]
\addplot[MemoBlue,thick] coordinates {
(2025-12-31,89.4981)
(2026-01-30,75.8448)
(2026-02-27,68.5199)
(2026-03-31,82.1465)
(2026-04-30,68.2904)
(2026-05-29,60.3529)
(2026-06-30,50.2294)
(2026-07-31,68.8302)
(2026-08-31,83.8631)
(2026-09-30,87.8732)
(2026-10-30,79.2654)
(2026-11-13,70.7976)
};
\addlegendentry{多空组合}
\addplot[orange,thick] coordinates {
(2025-12-31,102.5596)
(2026-01-30,108.3691)
(2026-02-27,110.3776)
(2026-03-31,99.8534)
(2026-04-30,110.6531)
(2026-05-29,117.9987)
(2026-06-30,118.7975)
(2026-07-31,109.1937)
(2026-08-31,113.8294)
(2026-09-30,107.0828)
(2026-10-30,94.0387)
(2026-11-13,98.1962)
};
\addlegendentry{等权基准}
\end{axis}
\end{tikzpicture}


\end{minipage}
\par\vfill

{\color{MemoRule}\hrule height 0.4pt}\medskip
\begin{minipage}[t]{0.36\textwidth}
\vspace{0pt}\small
\textbf{团队分工}\par\ReportAuthor\par\TeamName\par
\medskip\textbf{覆盖范围}\par\Coverage
\end{minipage}\hfill
\begin{minipage}[t]{0.60\textwidth}
\vspace{0pt}\footnotesize
\textbf{来源与日期。} 项目本地数据，由Kimi Datasource下载；\texttt{Data/CSI300\_daily/raw/}，12个CSV、24,634行，截至\DataDate。

\textbf{检查。} 日期错位、每月至少10只股票、样本外隔离、收益复算及有限数值检查均通过；不代表经济有效性通过。

\textbf{制备。} Python标准库计算；生成组研究、检查组复核、分配组汇总结论；AI辅助编程与报告排版。
\end{minipage}

\newpage
{\color{MemoBlue}\hrule height 0.5pt}\medskip
{\Large\bfseries 研究设计、候选比较与复现}\par\medskip

\textbf{因子定义}\par
20日反转因子为
\[
F_{i,t}=-\left(\frac{P_{i,t}}{P_{i,t-20}}-1\right),
\]
其中$P$为所给收盘价，$t$为月末，$t-20$为该股票此前第20个交易观测日。因子值越高，过去20日跌幅越大。每月末排序并分为五组，Q1最低、Q5最高；次月收益为下一个月末收盘价相对当月末收盘价的涨跌幅。停牌导致端点缺价时不进入当月检验。

\textbf{候选选择}\par
候选只根据样本内结果选择，最后12个月不参与选型。Rank IC是因子排名与次月收益排名的Spearman相关系数；ICIR为月均Rank IC除以月度Rank IC样本标准差，未年化。

\begin{center}
\captionof{table}{三个候选因子的样本内比较}\label{tab:candidates}
\begin{tabular}{@{}lrr@{}}
\toprule
候选因子 & 样本内月均Rank IC & 样本内ICIR \\
\midrule
20日动量 & -0.0149 & -0.0640 \\
\textbf{20日反转} & \textbf{0.0149} & \textbf{0.0640} \\
20日低波动 & -0.0307 & -0.1089 \\
\bottomrule
\end{tabular}

\end{center}

\textbf{样本外五分位表现}\par
各组按月等权并分别复利；多空组合则每月使用Q5收益减Q1收益后再复利，因此多空累计收益不等于Q5与Q1年化收益的简单差值。
\begin{center}
\captionof{figure}{样本外五分位组合年化收益}\label{fig:quintiles}
\begin{minipage}{0.72\textwidth}
\begin{tikzpicture}
\begin{axis}[ybar,width=\linewidth,height=48mm,bar width=12pt,symbolic x coords={Q1,Q2,Q3,Q4,Q5},xtick=data,ylabel={年化收益率（\%）},grid=major,grid style={gray!20},every axis label/.style={font=\small},tick label style={font=\small}]
\addplot[fill=MemoBlue,draw=none] coordinates {
(Q1,16.3524)
(Q2,-27.1594)
(Q3,-15.5329)
(Q4,3.0640)
(Q5,5.9738)
};
\end{axis}
\end{tikzpicture}

\end{minipage}
\end{center}
{\footnotesize 注：Q1至Q5年化收益依次为16.35\%、-27.16\%、-15.53\%、3.06\%、5.97\%，分组不单调。\par}

\textbf{风险指标定义}\par
年化收益按月度复利并以12个月年化；年化波动为月收益样本标准差乘$\sqrt{12}$；Sharpe使用零无风险利率；最大回撤基于月度多空净值。未扣成本，且A股实际做空受限，多空结果仅用于因子评价。

\textbf{复核与复现}\par
检查组确认所有35个月的收益日期晚于因子日期，每月均有34只有效股票；训练23个月、测试12个月互不重叠。月度结果见\texttt{output/monthly\_returns.csv}，完整数值及五项检查见\texttt{output/factor\_results.json}。在项目根目录运行\texttt{python src/factor\_research.py}可重算结果。

\textbf{后续建议}\par
优先补齐沪深300历史成分和明确复权的全样本数据，再检验行业中性、交易成本及不同持有期。当前结论应作为阴性研究结果保留，不建议继续调参直到样本外“变好看”。

\vfill
{\color{MemoRule}\hrule height 0.4pt}\medskip
{\footnotesize 数据来源：项目本地CSV；计算与图表：\texttt{src/factor\_research.py}。报告结果仅供课程研究，不构成投资建议。}
\end{document}
